\chapter{その他}
\section{歯車箱の厚さ}
歯車の厚さは、次の式で決定した。ここでCL=最終段中心距離=259.753[mm]となる。
\begin{eqnarray}
下部ケース&:&0.025CL + 3[mm] = 9.494 \approx 10[mm]\\
上部ケース&:&0.02CL + 3[mm] = 8.195 \approx 9[mm]
\end{eqnarray}
\section{歯車とケース内壁との最小間隔}
次の式で算出する。vは歯車収束である。
\begin{eqnarray}
第1段&:&C=2.5v+10[mm] =2.5 \times 6.7077 + 10 = 26.769[mm]  \approx 27[mm]\\
第2段&:&C=2.5v+10[mm] =2.5 \times 2.2113 + 10 = 15.528[mm] \approx 16[mm]
\end{eqnarray}
\section{歯車箱の放熱面積の決定}
\subsection{参考}
\begin{enumerate}
\item 馬力[HP],1[HP]=735.5[W]
\item 1[kcal/h]=1.163[W]
\item 1[inch]= 0.0254[m]
\item 1[mm]=0.03937[inch]
\end{enumerate}
\subsection{BS(British Standards)規格}
\begin{enumerate}
\item $\Delta t$ :許容温度と周囲温度の温度差
\begin{eqnarray}
\Delta t &=&  許容温度 - 周囲温度\\
         &=&  82℃ - 周囲温度
\end{eqnarray}
\item Q:歯車箱内での発熱量[kcal/h]
\begin{eqnarray}
Q=632(1-\eta)N
\end{eqnarray}
\item $\eta$:歯車装置の効率
\item N:歯車装置に与えられる馬力[HP]
\item A:歯車箱の放熱面積(底面を除く)[$m^2$]
\item K:熱通過係数($kcal/(m^2hK)$)
\end{enumerate}
BS規格では、放熱面積と歯車箱に加えられる馬力の間に次の関係がある。また、ここでは$\eta=0.98,\Delta t_{max} = 28,N=17000/735.5$とすると、次のようになる。
\begin{eqnarray}
A&=&\frac{Q}{K\Delta t_{max}}\\
 &=&\frac{632(1-\eta)N}{K\Delta t_{max}}\\
 &=&\frac{632\times(1-0.98) \times 17000/735.5}{10 \times 28}\\
 &=&1.0434[m^2]
\end{eqnarray}
\subsection{AGMA(American Gear Manufacturers Association)規格}
\begin{figure}[htbp]
  \begin{center}
\includegraphics[width=9cm]{../picture/20127161210319005.eps}
\end{center}
\caption{AGMA}
\end{figure}
AGMAの規格によれば、
\begin{eqnarray}
A=43.2{C_L}^{1.7}
\end{eqnarray}
である。ここで、$C_L(inch^2)=最終段中心距離$である。CL=最終段中心距離=259.753[mm]であるので、
\begin{eqnarray}
A=43.2 \times (259.753 \times 0.03937)^{1.7} &=& 2249.147 [inch^2]\\
&=&1.451[m^2]
\end{eqnarray}

\section{油面の高さの決定}
はねかけ式潤滑法では、油面の高さは中間軸の大歯車の最下位の歯丈の$2\sim3$倍程度にする.また、油面計をつける必要がある。
\section{重量計算}

\section{歯車箱への装着物}
\begin{enumerate}
\item 点検窓
\item 注油窓
\item 空気抜け(内圧上昇の防止,防塵防水に対する配慮)
\item 油面計
\item 排油口
\item 吊り金具
\item ノックピン(組み立て用)
\end{enumerate}
\section{仕上げ記号、はめあい記号の決定}

\section{参考文献}
\begin{enumerate}
\item \url{http://www.juntsu.co.jp/qa/qa2119.html}
\item \url{http://www.superior-inc.com/有限会社スピリアの構造変更情報館へようこそ！/構造変更一般/基本事項/強度検討書等を作成するための考察/圧縮（座屈）に付いて/}
\item \url{http://www.toishi.info/metal/hard_metal.html}
\item \url{http://kikakurui.com/b0/B0903-2001-01.html}
\end{enumerate}
